% Manual de Uso - Plataforma UBS
% Compilacao recomendada: XeLaTeX ou LuaLaTeX
\documentclass[12pt,a4paper]{article}

\usepackage[a4paper,margin=2.4cm]{geometry}
\usepackage{graphicx}
\usepackage{array}
\usepackage{enumitem}
\usepackage{titlesec}
\usepackage{fancyhdr}
\usepackage{tcolorbox}
\usepackage{hyperref}
\usepackage{tikz}
\usepackage[T1]{fontenc}
\usepackage[utf8]{inputenc}
\usepackage{lmodern}

% Tipografia (fallback para pdfLaTeX)
\newcommand{\headingfont}{\bfseries}

% Paleta (aproximada do frontend)
\definecolor{BrandBlue}{HTML}{2563EB}
\definecolor{BrandBlueDark}{HTML}{1D4ED8}
\definecolor{Ink}{HTML}{0F172A}
\definecolor{Slate}{HTML}{475569}
\definecolor{Paper}{HTML}{F8FAFC}
\definecolor{Line}{HTML}{E2E8F0}

\hypersetup{
  colorlinks=true,
  linkcolor=BrandBlueDark,
  urlcolor=BrandBlueDark,
  citecolor=BrandBlueDark
}

\pagestyle{fancy}
\fancyhf{}
\renewcommand{\headrulewidth}{0pt}
\fancyfoot[C]{\thepage}

\titleformat{\section}{\Large\headingfont\color{Ink}}{}{0pt}{}
\titleformat{\subsection}{\large\headingfont\color{Ink}}{}{0pt}{}

\newtcolorbox{FeatureBox}{
  colback=Paper,
  colframe=Line,
  boxrule=0.6pt,
  arc=8pt,
  left=10pt,
  right=10pt,
  top=8pt,
  bottom=8pt
}

\newcommand{\pill}[1]{\tikz[baseline]{\node[fill=BrandBlue!10,draw=BrandBlue!30,rounded corners=4pt,inner xsep=6pt,inner ysep=2pt, text=BrandBlueDark]{\small\textbf{#1}};}}

\begin{document}

% Capa
\begin{titlepage}
  \thispagestyle{empty}
  \begin{tikzpicture}[remember picture,overlay]
    \fill[BrandBlue] (current page.north west) rectangle ([yshift=-7cm]current page.north east);
    \fill[BrandBlueDark] (current page.north west) rectangle ([yshift=-0.6cm]current page.north east);
  \end{tikzpicture}

  \vspace*{2.6cm}
  {\headingfont\color{white}\fontsize{36}{42}\selectfont Manual de Uso\\}
  \vspace{0.4cm}
  {\headingfont\color{white}\fontsize{20}{24}\selectfont Plataforma UBS}

  \vspace{1.4cm}
  {\color{white}\large Versao atual \textemdash\ 07 de marco de 2026}

  \vfill
  {\color{white}\large Autoria: Caio e Samuel}

  \vspace*{1.2cm}
\end{titlepage}

\tableofcontents
\newpage

\section{Visao geral}
A Plataforma UBS oferece um ambiente completo para gestao de Unidades Basicas de Saude (UBS), integrando diagnosticos situacionais, agendamentos, comunicacao, materiais educativos e acompanhamento de cronogramas. O sistema foi pensado para diferentes perfis de uso, com fluxos controlados por permissao.

\section{Acesso e perfis}
\begin{FeatureBox}
\textbf{Perfis disponiveis}\par
\vspace{4pt}
\pill{USER} \pill{PROFISSIONAL} \pill{GESTOR} \pill{RECEPCAO} \pill{ACS}
\end{FeatureBox}

\subsection{Login}
\begin{itemize}[leftmargin=*]
  \item Acesse a tela de login e informe e-mail e senha.
  \item Em caso de esquecimento, utilize \textbf{Redefinir senha} no menu do usuario.
\end{itemize}

\subsection{Primeiro acesso}
\begin{itemize}[leftmargin=*]
  \item Usuarios recem-cadastrados iniciam com perfil \pill{USER}.
  \item Perfis de gestao devem ser concedidos por um gestor ou administrador.
\end{itemize}

\section{Navegacao principal}
A barra superior traz atalhos para as principais areas e o menu do usuario. O conteudo exibido varia conforme o perfil.

\begin{FeatureBox}
\textbf{Itens frequentes no menu}\par
\vspace{4pt}
\pill{Dashboard} \pill{Agendamento} \pill{Relatorios} \pill{Mapa de Problemas} \pill{Materiais} \pill{Cronograma} \pill{Suporte}
\end{FeatureBox}

\section{Funcionalidades por area}

\subsection{Dashboard}
Visao consolidada do desempenho da UBS, com indicadores e cards de acesso rapido.
\begin{itemize}[leftmargin=*]
  \item Acesse o \textbf{Dashboard} pelo menu superior.
  \item Revise os indicadores e clique nos cards para abrir atalhos.
  \item Use esta tela para definir prioridades do dia.
\end{itemize}

\subsection{Notificacoes \textemdash\ Gestor e Recepcao}
Centraliza avisos operacionais, solicitacoes e comunicados.
\begin{itemize}[leftmargin=*]
  \item Entre em \textbf{Notificacoes} no menu superior.
  \item Filtre por tipo de aviso e marque como lido apos resolver.
  \item Use para organizar a fila de demandas.
\end{itemize}

\subsection{Agendamento}
Permite criar, visualizar e organizar consultas e atendimentos.
\begin{itemize}[leftmargin=*]
  \item Acesse \textbf{Agendamento}.
  \item Clique em \textbf{Novo agendamento} e preencha paciente, data, horario e profissional.
  \item Salve e confirme o registro na lista ou no calendario.
  \item Para bloquear horarios, use \textbf{Bloquear horario} e selecione o periodo.
\end{itemize}

\subsection{Relatorios situacionais}
Gera relatorios com dados da UBS, anexos e metadados para acompanhamento tecnico.
\begin{itemize}[leftmargin=*]
  \item Acesse \textbf{Relatorios situacionais}.
  \item Selecione a UBS e o periodo de referencia.
  \item Preencha os campos e adicione anexos quando solicitado.
  \item Salve e utilize o relatorio para acompanhamento da gestao.
\end{itemize}

\subsection{Mapa de problemas e intervencoes}
Registro estruturado de problemas identificados e das intervencoes planejadas.
\begin{itemize}[leftmargin=*]
  \item Acesse \textbf{Mapa de problemas e intervencoes}.
  \item Cadastre o problema com descricao clara e local.
  \item Registre a intervencao planejada, responsavel e prazo.
  \item Atualize o status conforme a execucao evoluir.
\end{itemize}

\subsection{Materiais educativos}
Repositorio de materiais com anexos organizados por UBS.
\begin{itemize}[leftmargin=*]
  \item Acesse \textbf{Materiais educativos}.
  \item Clique em \textbf{Novo material} e adicione titulo, descricao e arquivos.
  \item Classifique por tema para facilitar a busca.
\end{itemize}

\subsection{Cronograma}
Planejamento de acoes, com etapas e prazos.
\begin{itemize}[leftmargin=*]
  \item Entre em \textbf{Cronograma}.
  \item Cadastre atividades, prazos e responsaveis.
  \item Atualize o andamento periodicamente para manter a equipe alinhada.
\end{itemize}

\subsection{Suporte e feedback}
Canal de comunicacao para duvidas, melhorias e suporte operacional.
\begin{itemize}[leftmargin=*]
  \item Acesse \textbf{Suporte e feedback}.
  \item Descreva o problema ou sugestao e envie.
  \item Acompanhe respostas quando disponiveis.
\end{itemize}

\subsection{Gestor de solicitacoes \textemdash\ Gestor}
Painel exclusivo para aprovar, recusar ou acompanhar solicitacoes administrativas.
\begin{itemize}[leftmargin=*]
  \item Acesse \textbf{Solicitacoes}.
  \item Analise o pedido e escolha \textbf{Aprovar} ou \textbf{Recusar}.
  \item Registre observacoes para historico de decisoes.
\end{itemize}

\subsection{Gerenciar mensagens \textemdash\ Recepcao}
Organiza mensagens recebidas e garante resposta rapida ao publico-alvo.
\begin{itemize}[leftmargin=*]
  \item Acesse \textbf{Mensagens}.
  \item Classifique por prioridade e responda conforme o protocolo local.
  \item Marque como resolvida ao concluir o atendimento.
\end{itemize}

\subsection{Gestao de equipes e microareas \textemdash\ Gestor e Recepcao}
Controle de equipes, vinculos com microareas e distribuicao de responsabilidades.
\begin{itemize}[leftmargin=*]
  \item Acesse \textbf{Gestao de equipes}.
  \item Cadastre profissionais e vincule a microarea correta.
  \item Revise periodicamente para manter o mapa atualizado.
\end{itemize}

\subsection{Setup da UBS \textemdash\ Profissional, Gestor e ACS}
Fluxo de configuracao inicial da UBS, necessario para liberar o uso da plataforma.
\begin{itemize}[leftmargin=*]
  \item Acesse \textbf{Setup da UBS} quando solicitado pelo sistema.
  \item Preencha identificacao da UBS e dados de contato.
  \item Salve a configuracao e valide no \textbf{Dashboard}.
\end{itemize}

\section{Fluxos do dia a dia (passo a passo)}

\subsection{Fluxo 1: Configuracao inicial da UBS}
\begin{itemize}[leftmargin=*]
  \item \textbf{Passo 1}: faca login com um perfil \pill{PROFISSIONAL}, \pill{GESTOR} ou \pill{ACS}.
  \item \textbf{Passo 2}: o sistema redireciona para \textbf{Setup da UBS} se ainda nao houver configuracao.
  \item \textbf{Passo 3}: preencha o nome da UBS, endereco, contatos e dados administrativos.
  \item \textbf{Passo 4}: revise os dados e clique em \textbf{Salvar}.
  \item \textbf{Passo 5}: volte ao \textbf{Dashboard} e confirme que os cards aparecem corretamente.
\end{itemize}
\textbf{Exemplo}: uma UBS recem-criada cadastra telefone e responsavel tecnico para liberar o acesso dos demais perfis.

\subsection{Fluxo 2: Criar um agendamento}
\begin{itemize}[leftmargin=*]
  \item \textbf{Passo 1}: acesse \textbf{Agendamento} no menu superior.
  \item \textbf{Passo 2}: clique em \textbf{Novo agendamento}.
  \item \textbf{Passo 3}: informe paciente, data, horario, profissional e tipo de atendimento.
  \item \textbf{Passo 4}: confirme os dados e salve.
  \item \textbf{Passo 5}: valide se o agendamento aparece na lista e no calendario.
\end{itemize}
\textbf{Exemplo}: atendimento de enfermagem agendado para 10/03, 09:30, com profissional Maria Souza.

\subsection{Fluxo 3: Bloquear horarios na agenda}
\begin{itemize}[leftmargin=*]
  \item \textbf{Passo 1}: em \textbf{Agendamento}, selecione \textbf{Bloquear horario}.
  \item \textbf{Passo 2}: defina data, hora inicial e hora final.
  \item \textbf{Passo 3}: informe o motivo do bloqueio (reuniao, treinamento, manutencao).
  \item \textbf{Passo 4}: salve e confirme o bloqueio no calendario.
\end{itemize}
\textbf{Exemplo}: bloqueio das 14:00 as 15:30 para reuniao de equipe.

\subsection{Fluxo 4: Registrar um relatorio situacional}
\begin{itemize}[leftmargin=*]
  \item \textbf{Passo 1}: acesse \textbf{Relatorios situacionais}.
  \item \textbf{Passo 2}: selecione a UBS e o periodo (ex.: fevereiro/2026).
  \item \textbf{Passo 3}: preencha os campos obrigatorios e descreva os principais indicadores.
  \item \textbf{Passo 4}: anexe documentos e evidencias, se necessario.
  \item \textbf{Passo 5}: salve e verifique o status do relatorio.
\end{itemize}
\textbf{Exemplo}: relatorio mensal com anexos de planilhas e fotos da unidade.

\subsection{Fluxo 5: Mapa de problemas e intervencoes}
\begin{itemize}[leftmargin=*]
  \item \textbf{Passo 1}: acesse \textbf{Mapa de problemas e intervencoes}.
  \item \textbf{Passo 2}: clique em \textbf{Novo problema} e descreva a situacao.
  \item \textbf{Passo 3}: informe local, impacto e prioridade.
  \item \textbf{Passo 4}: registre a intervencao com responsavel e prazo.
  \item \textbf{Passo 5}: atualize o status conforme o progresso.
\end{itemize}
\textbf{Exemplo}: problema de falta de vacinas com intervencao de reposicao em 7 dias.

\subsection{Fluxo 6: Publicar material educativo}
\begin{itemize}[leftmargin=*]
  \item \textbf{Passo 1}: acesse \textbf{Materiais educativos}.
  \item \textbf{Passo 2}: clique em \textbf{Novo material}.
  \item \textbf{Passo 3}: preencha titulo, descricao e tema.
  \item \textbf{Passo 4}: anexe arquivos (PDF, imagens ou documentos).
  \item \textbf{Passo 5}: salve e confirme que o material aparece na lista.
\end{itemize}
\textbf{Exemplo}: cartilha de prevencao da dengue para distribuicao na comunidade.

\subsection{Fluxo 7: Atualizar cronograma de acoes}
\begin{itemize}[leftmargin=*]
  \item \textbf{Passo 1}: acesse \textbf{Cronograma}.
  \item \textbf{Passo 2}: clique em \textbf{Nova atividade}.
  \item \textbf{Passo 3}: defina descricao, responsavel e prazo.
  \item \textbf{Passo 4}: salve e acompanhe o progresso.
  \item \textbf{Passo 5}: atualize o status semanalmente.
\end{itemize}
\textbf{Exemplo}: campanha de vacinacao com inicio em 15/03 e encerramento em 30/03.

\subsection{Fluxo 8: Resolver solicitacoes (Gestor)}
\begin{itemize}[leftmargin=*]
  \item \textbf{Passo 1}: acesse \textbf{Solicitacoes}.
  \item \textbf{Passo 2}: abra o pedido e revise os detalhes.
  \item \textbf{Passo 3}: selecione \textbf{Aprovar} ou \textbf{Recusar}.
  \item \textbf{Passo 4}: registre observacoes e salve.
  \item \textbf{Passo 5}: comunique a equipe sobre a decisao.
\end{itemize}
\textbf{Exemplo}: aprovacao de compra de materiais para sala de curativos.

\subsection{Fluxo 9: Atendimento de mensagens (Recepcao)}
\begin{itemize}[leftmargin=*]
  \item \textbf{Passo 1}: acesse \textbf{Mensagens}.
  \item \textbf{Passo 2}: leia a mensagem e categorize por prioridade.
  \item \textbf{Passo 3}: responda seguindo o protocolo local.
  \item \textbf{Passo 4}: marque como resolvida.
\end{itemize}
\textbf{Exemplo}: responder solicitacao de reagendamento enviada por um usuario.

\subsection{Fluxo 10: Enviar suporte e feedback}
\begin{itemize}[leftmargin=*]
  \item \textbf{Passo 1}: acesse \textbf{Suporte e feedback}.
  \item \textbf{Passo 2}: descreva o problema ou sugestao com detalhes.
  \item \textbf{Passo 3}: envie e aguarde retorno.
  \item \textbf{Passo 4}: acompanhe a resposta e confirme a solucao.
\end{itemize}
\textbf{Exemplo}: reportar erro ao anexar arquivos em relatorios.

\section{Boas praticas}
\begin{itemize}[leftmargin=*]
  \item Mantenha o cadastro da UBS atualizado para garantir relatorios coerentes.
  \item Revise o cronograma semanalmente com a equipe.
  \item Utilize materiais educativos padronizados para consistencia nas orientacoes.
  \item Registre problemas e intervencoes logo apos as visitas de campo.
\end{itemize}

\section{Seguranca e dados}
A plataforma utiliza autenticacao por token e controle de acesso baseado em perfil. Nunca compartilhe credenciais e encerre a sessao ao concluir o uso.

\section{Suporte}
Para suporte tecnico, utilize o modulo \textbf{Suporte e feedback}. Em emergencias operacionais, acione a gestao local.

\vfill
\begin{center}
  {\small \color{Slate}Documento interno da Plataforma UBS.}\par
\end{center}

\end{document}
